\chapter{Package cdtBook}
%!TEX root = ../docs.tex
\label{cg:book}

    El paquete cdtBook define un conjunto de comandos que permiten mostrar información acerca de la especificación de un proyecto, además muestra las personas/organizaciones encargadas de la elaboración del documento, revisión y aprobación del mismo; las fichas de proyecto y la elaboración del mismo.\\
    
\noindent Los paquetes requeridos para el funcionamiento adecuado de cdtBook son:
        
    \begin{itemize}
        \item {\bf longtable} Define el comando (environment) longtable, el cual es una versión de tabular multi-página
        \item {\bf colortbl} Permite utilizar colores en las tablas
        \item {\bf multirow} Provee una serie de extensiones al comando (environment) tabular como celdas multi-fila
        \item {\bf cdt/cdtTheme} Documentado en el capitulo \ref{ch:theme}
    \end{itemize}
    
\section{Datos de quien elabora}
    
    Permite especificar los datos generales de la organización en donde se elaboró/elaborá el documento de especificación de proyecto. Mediante el paquete se definen dos variables y un comando para asignarles valores (\refcmd{organizacion}).
    
\hypertarget{varOrganizacion}{}
\hypertarget{varCveOrganizacion}{}
%\subsection{varOrganizacion y varCveOrganizacion}
%\label{ssec:varOrg}

    \begin{itemize}
        \item {\bf\cmd{varOrganizacion}} Contiene el nombre completo de la organización.
        \item {\bf\cmd{varCveOrganizacion}} Contiene las siglas, abreviación o nombre corto de la organización.
    \end{itemize}
    
\subsection{organizacion}
\label{ssec:org}

    Este comando recibe dos argumentos, el primero (opcional) es la clave de la organización y el segundo especifica el nombre de la organización. El siguiente comando asignaría los valores 'Clave' y 'Nombre de la Organización' a las variables \refcmd{varCveOrganizacion} y \refcmd{varOrganizacion}, respectivamente.
    
    \begin{code}
\organizacion[CLAVE]{Nombre de la Organización}
    \end{code}
    
\noindent En caso de no especificarse el primer argumento, los valores asignados a \refcmd{varCveOrganizacion} y \refcmd{varOrganizacion} serían '' y 'Nombre de la Organización', respectivamente.
    
    \begin{code}
\organizacion{Nombre de la Organización}
    \end{code}
    
    
\section{Datos de quien elabora}
    
